\section{Trabalhos Relacionados}
A análise de dados aplicada à melhoria do desempenho acadêmico tem sido amplamente explorada na literatura, especialmente no contexto de Learning Analytics e Mineração de Dados Educacionais.

Na fase inicial deste trabalho, destacam-se como base teórica os estudos “Student Course Grade Prediction Using the Random Forest Algorithm: Analysis of Predictors' Importance” e “Analyzing Undergraduate Students' Performance Using Educational Data Mining”, que investigam o uso de técnicas de aprendizado de máquina para predição de desempenho acadêmico e identificação de fatores relevantes no progresso dos estudantes.

Essas abordagens fundamentam a análise e o tratamento dos dados dos discentes do curso de Ciência da Computação da UFRRJ, contribuindo para o treinamento do agente proposto. No entanto, diferentemente dos trabalhos mencionados, que se concentram predominantemente na predição de desempenho, este trabalho busca utilizar tais informações para apoiar a recomendação de disciplinas, com o objetivo de otimizar o tempo de formação dos estudantes.
